% ============================================================
\section{Faithful Observation and Classical-Limit Recovery}
\label{sec:classical-limit}
% ============================================================

The preceding chapters of Part~I developed the horizon calculus abstractly and
identified distinguished realized classes inside that grammar. The same
structural chain also exposed the native object of the calculus, the horizon
complex. A separate
question concerns the classical limit: when does a horizon system admit a
faithful continuum stage on which the observed laws recover the familiar
operator identities of standard calculus? This chapter answers that question in
two directions. It first shows how a classical operator model
induces a coherence-calculus description under finite observation. It then
proves that faithful towers recover the classical operator in the zero-defect
limit, and that directed horizon families admit canonical Hilbert completions
and regularity scales. The examples and comparison points used here lie in the
standard function-space and multiresolution literature
\cite{ReedSimon1972,AdamsFournier2003,Mallat1989,Daubechies1992}.

\subsection{Classical operators under finite observation}

The quickest route into the calculus from standard analysis is to begin with a
classical function space and impose a horizon cut.

\begin{example}[Classical derivative under subsampling]
\label{ex:classical-derivative-subsampling}
\lean{ClassicalDerivativeSubsamplingExample}
Let $\cH=L^2([0,1])$ and let $d=\frac{d}{dx}$ on a suitable Sobolev domain.
For each $h\in\bbN$, define $\CCPi{h}$ to be piecewise-constant averaging on
the dyadic partition into $2^h$ intervals. Then
\[
\CCObs{d}{h}=\CCRestr{d}{h}
\]
is the observed derivative at horizon $h$, while
\[
\CCDefect{d}{h}=\CCDefectExpr{d}{h}
\]
records exactly the failure of the observed derivative to close under that
finite observation cut \cite{AdamsFournier2003,Folland1999}.
\end{example}

\begin{example}[Nonlinear closure defect]
\label{ex:nonlinear-classical-defect}
\lean{NonlinearClassicalDefectExample}
Let $N$ be a nonlinear map on a function space, for example
$N(u)=u^2$. Projection need not commute with $N$, and the coherence defect
from Definition~\ref{def:coherence-defect} becomes
\[
\CCTau{h}{N}{u}:=\CCPi{h}N(u)-N(\CCPi{h}u).
\]
This is the exact closure term induced by finite observation. In the quadratic
case, Corollary~\ref{cor:quadratic-defect} identifies it with the within-block
variance created by the averaging projection.
\end{example}

\begin{remark}[Commuting filters and nonlinear obstruction]
\label{rem:commuting-filter-bridge}
\lean{commuting_filter_bridge}
If $\CCPi{h}$ is a commuting spectral filter for a linear operator $d$, then
the linear defect vanishes. The finite-observation obstruction becomes
nontrivial precisely when the observation map fails to commute with the
operator, or when the operator is nonlinear \cite{ReedSimon1972}.
\end{remark}

\subsection{Faithful observation and zero-defect recovery}

The reverse direction requires stronger hypotheses. A faithful tower is not
merely nested; it must recover the full state and the full operator graph in
the limit.

\begin{definition}[Faithful observation limit]
\label{def:faithful-observation-limit}
\lean{FaithfulObservationLimit}
A horizon tower $(\CCPi{h})_{h\in\bbN}$ on a Hilbert space $\cH$ is
\emph{faithful} if
\[
\CCPi{h}x\to x
\qquad\text{for every }x\in\cH.
\]
\end{definition}

\begin{proposition}[Faithful recovery of a classical operator]
\label{prop:faithful-operator-recovery}
\lean{faithful_operator_recovery}
Let $d:\mathcal D(d)\subseteq \cH\to \cH$ be a closed linear operator on a
Hilbert space, and let $(\CCPi{h})_{h\in\bbN}$ be orthogonal projections such
that:
\begin{enumerate}[label=(\roman*), leftmargin=2em, itemsep=0.2em]
\item the tower is faithful in the sense of
      Definition~\ref{def:faithful-observation-limit};
\item $\CCPi{h}\mathcal D(d)\subseteq \mathcal D(d)$ for every $h$;
\item for every $x\in\mathcal D(d)$,
      \[
      d\,\CCPi{h}x\to dx
      \qquad\text{in }\cH.
      \]
\end{enumerate}
Then for every $x\in\mathcal D(d)$:
\[
\|\CCPibar{h}d\CCPi{h}x\|\to 0,
\qquad
\CCObs{d}{h}x=\CCPi{h}d\CCPi{h}x\to dx.
\]
In particular, the projection-induced defect vanishes on the graph of $d$ in
the faithful limit.
\end{proposition}

(Proof in Appendix~\ref{sec:appendix-continuum}.)

\begin{remark}[Faithful limit of returned residual]
\label{rem:faithful-limit-returned-residual}
\lean{faithful_limit_returned_residual}
Proposition~\ref{prop:faithful-operator-recovery} is the bridge-side reason
that a residual return field is not a permanent extra source sector. At finite
observation, leakage into the shadow and its returned visible trace may be
nonzero. But on the graph of the recovered operator, faithfulness forces the
projection-induced defect to vanish in the limit. Thus any residual return
field attached to the finite-horizon observed law is a genuine eliminated datum
whose unrecovered part disappears along a faithful bridge, rather than an
independent constitutive forcing term.
\end{remark}

\begin{remark}[External parameters remain optional]
\label{rem:external-parameter-optional}
\lean{ExternalParameterIdentification}
No physical time parameter is forced here. If a classical model already carries
an external parameter $t$, then Coherence Calculus may be applied in another
direction, for example spatial coarse-graining or spectral truncation. Conversely,
an ordinal refinement index $n$ may be identified with an external parameter by
choosing a step size and writing $t=n\Delta t$. Such an identification is
external to the calculus itself.
\end{remark}

\subsection{Directed horizon families and Hilbert completion}

Faithful towers arise naturally from directed systems of finite horizons.
Passing to the limit produces a canonical Hilbert stage together with an exact
tail-energy budget.

\begin{definition}[Directed horizon family]
\label{def:directed-horizon-family}
\lean{DirectedHorizonFamily}
A \emph{directed horizon family} consists of:
\begin{enumerate}[label=(\roman*), leftmargin=2em, itemsep=0.2em]
\item a sequence of finite-dimensional inner-product spaces
      $(\CCHobs{h})_{h\in\bbN}$;
\item isometric inclusions
      \[
      \iota_{h\to h+1}:\CCHobs{h}\hookrightarrow \CCHobs{h+1}
      \qquad(h\in\bbN).
      \]
\end{enumerate}
Its algebraic direct limit is
\[
\cH_{\mathrm{alg}}:=\bigcup_{h\in\bbN}\CCHobs{h},
\]
where each stage is identified with its image under the inclusion maps.
\end{definition}

\begin{theorem}[Hilbert completion of a directed horizon family]
\label{thm:hilbert-completion}
\lean{hilbert_completion_of_directed_horizon_family}
Let $(\CCHobs{h},\iota_{h\to h+1})$ be a directed horizon family.
\begin{enumerate}[label=(\roman*), leftmargin=2em, itemsep=0.2em]
\item The completion
      \[
      \cH_\infty:=\overline{\cH_{\mathrm{alg}}}
      \]
      is a separable Hilbert space.
\item Each stage $\CCHobs{h}$ embeds isometrically into $\cH_\infty$.
\item Orthogonal projection onto the stage subspace defines maps
      \[
      \CCPi{h}:\cH_\infty\to \CCHobs{h}
      \]
      forming a transitive horizon tower on $\cH_\infty$.
\item The resulting tower is faithful:
      \[
      \CCPi{h}x\to x
      \qquad(x\in\cH_\infty).
      \]
\end{enumerate}
\end{theorem}

(Proof in Appendix~\ref{sec:appendix-continuum}.)

\begin{corollary}[Exact horizon error budget]
\label{cor:horizon-error-budget}
\lean{exact_horizon_error_budget}
Let $(\CCPi{h})$ be the faithful tower produced by
Theorem~\ref{thm:hilbert-completion}, and define increment projectors by
\[
\CCDelta{h}:=\CCPi{h}-\CCPi{h-1},
\qquad
\CCPi{-1}:=0.
\]
Then for every $x\in\cH_\infty$ and every initial horizon $h_0$,
\[
\|x-\CCPi{h_0}x\|^2
=
\sum_{j\ge h_0}\|\CCDelta{j+1}x\|^2.
\]
So the full tail energy is exactly the sum of the layer energies.
\end{corollary}

\begin{proof}
Apply the orthogonal telescoping identity from
Lemma~\ref{lem:orthogonal-telescoping} to the faithful tower on $\cH_\infty$.
The refinement bands are orthogonal, so the tail norm splits as the sum of the
squared increment norms.
\end{proof}

\subsection{Tower-induced regularity ladders}

The exact tail decomposition defines a natural smoothness scale without
appealing to coordinates or derivatives on the ambient space.

\begin{definition}[Tower regularity space]
\label{def:tower-regularity}
\lean{TowerRegularitySpace}
Let $\cH_\infty$ be the completion from
Theorem~\ref{thm:hilbert-completion}. For $s\ge 0$ and a weight function
$w_s:\bbN\to(0,\infty)$ with $w_s(h)\to\infty$ as $h\to\infty$, define
\[
\cH^{(s)}
:=
\left\{
x\in\cH_\infty:
\|x\|_{(s)}^2
:=
\sum_{h\ge 0} w_s(h)\,\|\CCDelta{h}x\|^2
<
\infty
\right\}.
\]
Typical choices are:
\begin{itemize}[itemsep=0.2em]
\item dyadic weights $w_s(h)=2^{2sh}$ for dyadically refined towers;
\item polynomial weights $w_s(h)=(1+h)^{2s}$ for generic towers.
\end{itemize}
\end{definition}

\begin{proposition}[Regularity ladder properties]
\label{prop:regularity-ladder}
\lean{regularity_ladder_properties}
Under Definition~\ref{def:tower-regularity}:
\begin{enumerate}[label=(\roman*), leftmargin=2em, itemsep=0.2em]
\item $\cH^{(s)}$ is a Hilbert space with inner product
      \[
      \langle x,y\rangle_{(s)}
      :=
      \sum_{h\ge 0} w_s(h)\,\langle \CCDelta{h}x,\CCDelta{h}y\rangle;
      \]
\item if $w_0(h)\equiv 1$, then $\cH^{(0)}=\cH_\infty$;
\item if $s'>s$ and $w_{s'}(h)\ge w_s(h)$ for all $h$, then
      \[
      \cH^{(s')}\subseteq \cH^{(s)}
      \]
      continuously;
\item the intersection
      \[
      \cH^{(\infty)}:=\bigcap_{s\ge 0}\cH^{(s)}
      \]
      consists exactly of those elements whose increment norms decay faster than
      every chosen weight scale.
\end{enumerate}
\end{proposition}

\begin{proof}
The increment decomposition identifies each $\cH^{(s)}$ with a weighted
$\ell^2$ space of refinement bands. Completeness and the stated embeddings are
the standard properties of weighted $\ell^2$ norms.
\end{proof}

\begin{remark}[Relation to classical smoothness scales]
\label{rem:regularity-model}
\lean{RegularityModel}
In standard multiresolution settings, weighted increment norms recover familiar
regularity scales such as Sobolev and Besov spaces. Here the same ladder is
defined purely from the horizon tower and its exact tail decomposition, without
putting classical derivatives into the axioms
\cite{AdamsFournier2003,Triebel1983,Mallat1989,Daubechies1992}.
\end{remark}

\subsection{A standard dyadic model}

The abstract continuation schema has a canonical model on a classical function
space.

\begin{definition}[Dyadic tower on \texorpdfstring{$L^2([0,1]^n)$}{L2([0,1]^n)}]
\label{def:dyadic-l2-tower}
\lean{DyadicL2Tower}
Let $\cH=L^2([0,1]^n)$ for $n\ge 1$. For each $h\in\bbN$, partition $[0,1]^n$
into $2^{nh}$ dyadic cubes of side length $2^{-h}$ and define $\CCPi{h}$ to be
the orthogonal projection onto the piecewise-constant functions on that
partition:
\[
(\CCPi{h}f)(x)
:=
\frac{1}{|Q|}\int_Q f\,d\mu,
\]
where $Q$ is the unique dyadic cube containing $x$.
\end{definition}

\begin{proposition}[Standard properties of the dyadic model]
\label{prop:dyadic-tower-properties}
\lean{dyadic_tower_properties}
The dyadic tower from Definition~\ref{def:dyadic-l2-tower} satisfies:
\begin{enumerate}[label=(\roman*), leftmargin=2em, itemsep=0.2em]
\item each $\CCPi{h}$ is an orthogonal projection with
      \[
      \dim(\Img(\CCPi{h}))=2^{nh};
      \]
\item transitivity:
      \[
      \CCPi{h}\CCPi{h'}=\CCPi{\min(h,h')};
      \]
\item faithfulness:
      \[
      \CCPi{h}f\to f
      \qquad\text{strongly in }L^2([0,1]^n).
      \]
\end{enumerate}
Hence the dyadic tower realizes the abstract continuation schema of
Theorem~\ref{thm:hilbert-completion}.
\end{proposition}

\begin{proof}
Orthogonal projection onto piecewise-constant functions gives (i), nesting of
the dyadic partitions gives (ii), and the Lebesgue differentiation theorem
implies (iii) \cite{Folland1999}.
\end{proof}

\begin{remark}[Model status]
\label{rem:dyadic-model-status}
\lean{DyadicModelStatus}
This dyadic tower is a model of the Part~II interface, not a derivation from the
bedrock axioms of Part~I. Its role is to show concretely how the abstract
horizon grammar recovers a familiar classical function-space setting once a
faithful tower is supplied.
\end{remark}
